\section{Related Literature}

This paper contributes to several interconnected literatures spanning fertility, housing markets, and political economy. We organize the discussion around five strands and clarify how our framework unifies them.

\medskip
\noindent\textbf{Fertility and housing costs.}
A growing body of evidence documents the link between housing costs and fertility outcomes. Ermisch (1988) provided early evidence that housing costs affect the timing of childbearing in the United Kingdom, showing that couples facing higher housing expenditure delay family formation. Simon and Tamura (2009) extended this line of inquiry to U.S.\ metropolitan areas, finding that local housing costs are negatively associated with fertility rates across cities. Dettling and Kearney (2014) sharpened the identification using MSA-level variation in house prices and established that rising house prices reduce birth rates among non-homeowners while modestly increasing them among owners, consistent with a wealth-effect channel operating alongside the cost-of-space channel. Clark (2012) documented similar patterns in the UK, emphasizing that binding housing constraints---particularly difficulty accessing owner-occupied housing---delay partnership formation and the transition to parenthood.

Lovenheim and Mumford (2013) exploited household-level variation in housing wealth from the PSID and found that a \$100,000 increase in home equity raises the probability of having a child by 16--18~percent among homeowners, with no comparable effect for renters, confirming that tenure status mediates the fertility--house-price relationship. Daysal, Lovenheim, Siersbæk, and Wasser (2021) replicated this asymmetry in Danish register data and showed that housing wealth gains also improve infant health at birth, suggesting that the channel operates partly through resource availability rather than pure opportunity cost. These findings sit alongside a broader literature connecting local economic conditions to demographic outcomes. Autor, Dorn, and Hanson (2019) showed that adverse labor market shocks reduce marriage and fertility, highlighting how economic insecurity transmits into household formation decisions. Doepke and Kindermann (2019) developed a quantitative model in which disagreement between partners over the costs of childrearing contributes to fertility decline, underscoring that the decision to have children is shaped by both partners' perceptions of economic burden. Kearney, Levine, and Pardue (2022) surveyed the post-2007 U.S.\ birth-rate decline and concluded that no single economic or policy change accounts for the sustained fall; their analysis leaves open the possibility that persistent structural forces---such as housing scarcity---cumulate beneath the cohort-preference explanation.

Our paper builds directly on this strand by asking how a specific, persistent source of housing cost pressure---supply restrictions driven by political incumbents---affects not just the level of fertility but the timing of first births.

\medskip
\noindent\textbf{Birth timing and completed fertility.}
Our mechanism operates through delayed first births: when housing is expensive, households postpone the first child, and the finite reproductive window means some postponed births become foregone births. This builds on the quantity--quality framework of Becker (1960) and Becker and Lewis (1973), in which the shadow price of children rises with input costs. Hotz, Klerman, and Willis (1997) emphasized the role of tempo effects---delays in childbearing compress the remaining window for subsequent births and reduce completed family size. Goldin and Katz (2002) showed that access to the contraceptive pill enabled women to invest in careers, contributing to delayed family formation, while Caucutt, Guner, and Knowles (2002) showed that matching frictions lower fertility through a similar delay channel; both demonstrated that delay translates into lower completed fertility.

Several quantitative macro models have explored the aggregate implications of fertility choice. Greenwood, Seshadri, and Vandenbroucke (2005) explained the baby boom and bust through changes in the cost of household production, demonstrating that technology shocks that reduce the time cost of children can generate large swings in aggregate fertility---a mechanism that operates through the same shadow-price channel our model exploits, but on the technology side rather than the housing side. De la Croix and Doepke (2003) showed that differential fertility across income groups is a quantitatively important channel linking inequality to growth, establishing that models with endogenous fertility can generate macroeconomically significant effects even when individual-level responses appear small. Jones and Schoonbroodt (2016) embedded fertility in a dynastic model with aggregate shocks and showed that the Great Depression alone can generate a baby bust of 58~percent of the observed U.S.\ magnitude, followed by a substantial boom---confirming that fertility responds meaningfully to persistent economic conditions in calibrated general-equilibrium settings. Our contribution is to identify housing scarcity as an additional source of delay and to trace its aggregate fertility consequences in a political-economy framework where the source of scarcity is itself endogenous.

\medskip
\noindent\textbf{Housing supply regulation and NIMBYism.}
The source of housing scarcity in our framework is endogenous supply restrictions, a mechanism grounded in the housing supply regulation literature. Glaeser and Gyourko (2003) documented that land-use regulations raise housing costs substantially above construction costs in many U.S.\ cities. Glaeser, Gyourko, and Saks (2005) showed that the high cost of housing in Manhattan is attributable primarily to regulatory constraints rather than land scarcity or construction costs, and that these constraints have tightened over time. Saiz (2010) decomposed the housing supply elasticity into geographic and regulatory components, finding that regulation is the binding constraint in the most expensive markets. Gyourko, Hartley, and Krimmel (2021) updated and extended the Wharton Residential Land Use Regulatory Index to nearly 2,500 communities, documenting that the regulatory environment has tightened further since the mid-2000s and that the most regulated markets are disproportionately high-cost coastal cities. Turner, Haughwout, and van der Klaauw (2014) decomposed the welfare effects of land-use regulation into own-lot, external, and supply components, finding that the supply restriction effect generates large welfare losses---a result that motivates our focus on the quantity margin of housing supply rather than amenity externalities. The aggregate consequences of local supply restrictions are potentially large: Hsieh and Moretti (2019) estimated that housing constraints in high-productivity cities reduce U.S.\ aggregate output by misallocating labor across space.

Our paper takes the political-economy model of NIMBYism developed in Gross and Chivers (2025) as its upstream benchmark. That framework showed how majority-homeowner communities vote to restrict new construction in order to protect incumbents' house values, generating a wedge between efficient and equilibrium housing supply. We do not revisit that political-economy mechanism in detail here; rather, we embed it as the supply side of a richer demographic model and ask what happens to fertility when the housing stock is determined by incumbent voting.

\medskip
\noindent\textbf{Political economy of housing.}
The idea that homeowners have incentives to restrict supply has deep roots. Fischel (2001) articulated the ``homevoter hypothesis,'' arguing that homeowners treat their house as their primary financial asset and vote on local land-use policy accordingly, favoring restrictions that protect property values. Ortalo-Magn\'{e} and Prat (2014) formalized the political economy of urban growth, showing how the voting power of incumbent owners can generate inefficiently low levels of housing construction. Hankinson (2018) provided the first experimental evidence on NIMBYism, showing that even renters in high-cost cities oppose nearby construction at rates comparable to homeowners---evidence of scale-dependent preferences that help explain why supply restrictions persist even when a majority of residents support more housing in the abstract. Marble and Nall (2021) linked deed-level homeownership records to voter turnout for 18~million individuals and found that purchasing a home substantially increases participation in local elections, with the turnout boost nearly doubling when zoning issues are on the ballot, suggesting that homeowners' political engagement is partly asset-motivated. Parkhomenko (2023) embedded these mechanisms in a spatial equilibrium framework with heterogeneous agents, demonstrating that politically determined supply restrictions redistribute welfare from renters to owners and reduce aggregate housing consumption.

Our model inherits this political-economy structure but extends it to a setting where the demographic composition of the electorate---which depends on fertility outcomes---feeds back into housing policy. When supply restrictions delay childbearing, the resulting smaller cohorts alter the future balance of owners and renters, potentially reinforcing or attenuating the original supply constraint.

\medskip
\noindent\textbf{Demographic feedback into housing markets.}
Mankiw and Weil (1989) first argued that the baby-boom cohort's entry into the housing market drove the 1970s price boom, and that the subsequent bust would follow as smaller cohorts matured---a prediction that proved directionally correct even if its timing was off. Green and Hendershott (1996) refined the age--demand relationship, showing that the age profile of housing demand peaks in middle age and that shifts in cohort size generate persistent price effects. Chiuri and Jappelli (2003) documented that mortgage-market development shifts the age profile of homeownership toward the young across 14~OECD countries, implying that financial frictions interact with demographics to shape housing demand. Sommer and Sullivan (2018) showed that eliminating the U.S.\ mortgage interest deduction would lower homeownership rates and house prices substantially, providing a quantitative benchmark for the magnitude of policy-driven price effects. Mumtaz and Sust\v{e}k (2024) developed a model in which demographic structure is a first-order driver of house-price dynamics.

These papers treat demographics as exogenous. Our framework closes the loop: housing scarcity delays first births and reduces cohort size, which alters future housing demand and the political equilibrium that determines supply. The feedback is quantitatively modest in our calibration, but the mechanism is novel and correctly signed.

\medskip
\noindent\textbf{Contribution.}
To our knowledge, this is the first paper to embed an endogenous fertility decision---with explicit first-birth timing---into a political-economy model of housing supply with endogenous land-use restrictions. The fertility--housing-cost literature (strand one) has documented the empirical link between house prices and birth rates but has not modeled the supply-side origin of high prices. The birth-timing literature (strand two) has shown that delay reduces completed fertility but has not connected delay to a specific, policy-relevant source of housing cost pressure. The housing supply regulation and political-economy literatures (strands three and four) have modeled endogenous supply restrictions and their welfare consequences but have treated household size as fixed. The demographic feedback literature (strand five) has shown that cohort size drives housing demand but has taken demographics as exogenous. By connecting these strands, we can study a channel that none of them addresses in isolation: NIMBYism raises housing costs, which delays first births and reduces completed fertility, which reshapes the future electorate and housing demand. The quantitative framework allows us to assess the magnitude of this channel and to evaluate policies that target the supply side of the housing market as instruments of demographic policy.
